\section*{Chapitre \ref{ch:b2:batteries-electrolysis} --- Piles, piles à combustible et électrolyse}

\begin{solution}{exo:b2:batteries-electrolysis:1}
Zinc seul~: environ \qty{-0.82}{V} et un courant de 0.15 (unités du modèle). Au contact du
cuivre~: environ \qty{-0.74}{V} et 2.4. Le courant, donc la vitesse de dissolution, est
environ seize fois plus grand au contact du cuivre.
\end{solution}

\begin{solution}{exo:b2:batteries-electrolysis:2}
$m = 500 \times 3600 \times 65.38/(2 \times 96\,485) = \qty{610}{g}$.
\end{solution}

\begin{solution}{exo:b2:batteries-electrolysis:3}
$E^\circ = \qty{2.04}{V}$ (comme dans le chapitre)~; six éléments en série donnent
\qty{12}{V}.
\end{solution}

\begin{solution}{exo:b2:batteries-electrolysis:4}
$1/6.94 = \qty{0.144}{mol}$~; $0.144 \times 96\,485 = \qty{1.39e4}{C}$, soit
\qty{3.86}{A.h}.
\end{solution}

\begin{solution}{exo:b2:batteries-electrolysis:5}
À 100~\%~: $200 \times 86\,400 \times 63.546/(2 \times 96\,485) = \qty{5.69}{kg}$~;
$\eta_F = 5.50/5.69 = 0.967$.
\end{solution}

\begin{solution}{exo:b2:batteries-electrolysis:6}
$w = 2 \times 96\,485 \times 3.3/(0.90 \times 0.06538) = \qty{1.08e7}{J/kg}$, soit
\qty{3.0}{kW.h/kg}.
\end{solution}

\begin{solution}{exo:b2:batteries-electrolysis:7}
$U = E - rI$~: $r = (1.55 - 1.40)/(0.60 - 0.10) = \qty{0.30}{\ohm}$~; $E = 1.55 + 0.30
\times 0.10 = \qty{1.58}{V}$.
\end{solution}

\begin{solution}{exo:b2:batteries-electrolysis:8}
Anode \ce{2Cl- -> Cl2 + 2e-}~; cathode \ce{2H2O + 2e- -> H2 + 2OH-}. $\Delta_r G^\circ =
2(-157.24) - 2(-131.23) - 2(-237.13) = \qty{422.2}{kJ/mol}$, $U_{\min} =
\qty{2.19}{V}$. Le dichlore réagirait avec l'hydroxyde (\ce{Cl2 + 2OH- -> Cl- +
ClO- + H2O}), ce qui gâcherait les deux produits, et un mélange de dichlore et de dihydrogène
peut exploser.
\end{solution}

\begin{solution}{exo:b2:batteries-electrolysis:9}
\ce{NaCl -> Na + 1/2Cl2}, $n = 1$~: $U_{\min} = 310\,674/96\,485 = \qty{3.22}{V}$.
Par kilogramme de sodium~: $96\,485 \times 3.22/0.02299 = \qty{1.35e7}{J}$, soit
\qty{3.75}{kW.h}.
\end{solution}

\begin{solution}{exo:b2:batteries-electrolysis:10}
Sur une cathode de zinc, la réduction de \ce{H+} est lente et son mur se situe au-dessous de la
vague du zinc~: quand on abaisse le potentiel, le zinc se dépose en premier, avec peu de
dihydrogène. Sur platine, la réduction de \ce{H+} est rapide et commence vers
\qty{0}{V}~: au début, seul du dihydrogène se dégage. Une fois le platine recouvert de
zinc, la cathode se comporte comme une cathode de zinc et le zinc se dépose.
\end{solution}

\begin{solution}{exo:b2:batteries-electrolysis:11}
À \qty{1250}{K}~: $\Delta_f G^\circ(\ce{Al2O3}) = -1328.3 + 0.75 \times 66.0 =
\qty{-1278.8}{kJ/mol}$, $\Delta_f G^\circ(\ce{CO2}) = \qty{-396.1}{kJ/mol}$. $\Delta_r G^\circ =
3(-396.1) - 2(-1278.8) = \qty{1369}{kJ}$, $n = 12$~: $U_{\min} = \qty{1.18}{V}$. Anode inerte~:
$\Delta_r G^\circ = \qty{2557.5}{kJ}$, $U_{\min} = \qty{2.21}{V}$~; la combustion du carbone
fournit environ \qty{1.03}{V} du travail.
\end{solution}

\begin{solution}{exo:b2:batteries-electrolysis:12}
Par mole d'eau, la charge $2F$ est la même dans les deux sens~: rendement $0.70/1.80 =
0.39$. Le reste est de la chaleur~: les \omterm{def:b2:current-potential-curves:overpotential}{surtensions} des deux électrodes à oxygène (lentes dans
les deux sens), celles des électrodes à hydrogène, et les pertes ohmiques.
\end{solution}

\begin{solution}{pb:b2:batteries-electrolysis:1}
\textbf{1.} Aluminium $+3 \to 0$~; l'oxygène reste à $-2$~; carbone $0 \to +4$.
\textbf{2.} Cathode \ce{Al^3+ + 3e- -> Al}~; anode \ce{C + 2O^2- -> CO2 + 4e-}.
\textbf{3.} \ce{2Al2O3 + 3C -> 4Al + 3CO2}, $n = 12$.
\textbf{4.} \ce{Al2O3}~: $-1328.3 + 0.75 \times (1328.3 - 1262.3) = \qty{-1278.8}{kJ/mol}$~;
\ce{CO2}~: \qty{-396.1}{kJ/mol}.
\textbf{5.} $3(-396.1) - 2(-1278.8) = \qty{+1369}{kJ}$.
\textbf{6.} $1\,369\,000/(12 \times 96\,485) = \qty{1.18}{V}$.
\textbf{7.} $2 \times 1278.8/(12 \times 96.485) = \qty{2.21}{V}$~: l'oxydation du carbone,
elle-même \omterm{def:b2:reaction-free-energy:exergonic}{exergonique}, paie près de la moitié du travail de décomposition de l'oxyde.
\textbf{8.} $10^6/26.982 = \qty{3.706e4}{mol}$.
\textbf{9.} $3 \times 96\,485 \times 3.706 \times 10^4 = \qty{1.073e10}{C}$.
\textbf{10.} $1.073 \times 10^{10}/0.94 = \qty{1.141e10}{C}$.
\textbf{11.} $1.141 \times 10^{10}/3.0 \times 10^5 = \qty{3.80e4}{s}$, \qty{10.6}{h}.
\textbf{12.} $24/10.6 = 2.27$~: \qty{2.27}{t} par jour.
\textbf{13.} Une partie de l'aluminium dissous dans le bain atteint le gaz anodique et
est réoxydée par le dioxyde de carbone~: sa charge est donc dépensée deux fois~; une partie du courant
fuit aussi à travers le bain sans l'électrolyser.
\textbf{14.} $1.141 \times 10^{10} \times 4.1 = \qty{4.68e10}{J}$, \qty{13.0}{MW.h}.
\textbf{15.} \qty{13.0}{kW.h/kg}.
\textbf{16.} $3 \times 96\,485 \times 1.18/0.026982 = \qty{1.27e7}{J/kg}$, \qty{3.5}{kW.h/kg}.
\textbf{17.} En chaleur~: les \omterm{def:b2:current-potential-curves:overpotential}{surtensions} des électrodes et, surtout, la
chute ohmique dans le bain. Cette chaleur maintient le bain fondu vers
\qty{1250}{K}, ce qui exigerait sinon du combustible.
\textbf{18.} $4.1 \times 3.0 \times 10^5 = \qty{1.23}{MW}$.
\textbf{19.} $28\,690/0.026982 = \qty{1.06e6}{J/kg}$, \qty{0.30}{kW.h/kg}~: environ 2~\% de
l'énergie de l'\omterm{def:b2:batteries-electrolysis:electrolysis}{électrolyse}.
\textbf{20.} $1.18/4.1 = 0.29$.
\textbf{21.} $3.706 \times 10^4 \times 3/4 \times 12.011 = \qty{334}{kg}$.
\textbf{22.} $3.706 \times 10^4 \times 3/4 \times 44.009 = \qty{1.22e3}{kg}$.
\textbf{23.} Le haut des anodes, chaud, brûle dans l'air~; une partie du carbone part
sous forme de monoxyde de carbone (\ce{C + CO2 -> 2CO} à cette température), ce qui consomme deux fois
plus de carbone par atome d'oxygène.
\textbf{24.} \qty{1.22}{kg} de \ce{CO2} par kilogramme d'aluminium.
\textbf{25.} La tension minimale augmenterait d'environ \qty{1.03}{V}, mais l'anode
dégagerait du dioxygène au lieu de dioxyde de carbone et ne serait pas consommée.
\textbf{26.} La cuve consomme $\boldsymbol{\approx \qty{13}{kW.h}}$ d'énergie électrique
par kilogramme d'aluminium.
\end{solution}
